\section{A uniform bound for the number of separating polygons}
\label{sec:06-all-count}

The fixed-fugacity cylinder asymptotic determines the nesting exponent.
A different consequence of the same finite sector identity controls the
number of separating polygons, including counts that grow with the
width. For that purpose the fugacity must be allowed to grow without a
fixed upper bound. The real, untruncated sector estimate provides this
uniformity directly.

Use the physical balanced two-marker cylinder with even width
\(N=2m\), band parameters \(2\lambda\) on the first half and
\(\lambda\) on the second, and \(\lambda=\pi/8\). For each integer
\(a\ge0\), let \(w_N(a)\) be the critical mass of families of exactly
\(a\) pairwise vertex-disjoint polygons, each separating the two
markers. The weight of a family is \(\prod_P\rho^{|P|}\), with
\(\rho=(2+\sqrt2)^{-1/2}\) per polygon vertex. In particular the
factor two for a polygon is not included in \(w_N(a)\). The physical
partition is
\begin{equation}
 Z_N(\chi)=\sum_{a\ge0}w_N(a)\chi^a,
 \qquad w_N(a)\ge0,
 \qquad w_N(0)=1.
 \label{eq:06-positive-coefficients}
\end{equation}
This is the separating-loop contraction defined in
Section~\ref{sec:02-cylinder}. At fixed \(N\) it is a polynomial with
finite coefficients; a cut between the two markers bounds the number
of disjoint separating polygons.

\begin{theorem}[Uniform fugacity bound]\label{E:all-fugacity-bound}
There are constants \(C<\infty\) and \(N_0\) such that, for every
even \(N\ge N_0\) and every real \(v\ge0\),
\begin{equation}
 Z_N(2\cosh v)
 \le \exp\!\bigl(C\log N\,(1+v^2)\bigr).
 \label{eq:06-all-fugacity}
\end{equation}
The constants are independent of \(v\).
\end{theorem}
\begin{proof}
Set \(L=(\log N)/\gamma\), where \(\gamma=4/3\), and fix the
smoothing parameter used in the real-sector bound. It is fixed
independently of \(N\) and \(v\). We use the exact sector identity
\eqref{eq:G-22} at \(\eta=iv\), retaining all valid integer sectors
and dimensions. The imbalance \(M\) remains a real integer throughout
this proof. Thus \(v=\operatorname{Im}\eta\), whereas
\(\operatorname{Im}M=0\).

Use the normalization \eqref{eq:04-physical-normalization},
\[
 G_N(\eta)=\Delta_N Z_N(2\cos\eta).
\]
The scalar \(\Delta_N\) is nonzero and independent of \(\eta\).
We first bound \(G_N(iv)\), and then divide by \(|\Delta_N|\).

Recall the quantities in that identity. The explicit configuration
\(\zeta\) consists of exterior main particles and all far particles;
\(n_e\) counts its integration variables, with a fused double counted
as one variable, and
\[
 \ell_e=\sum_{z\in\zeta}(|\operatorname{Re}z|-L)_+.
\]
The interior dimensions \(m'_\sigma(M)\) are nonnegative integers
at most \(N\). Each explicit species retains its factorial divisor;
its elementary multiplicity is one for a singleton and two for a
fused double. The frozen prefactor has modulus at most
\(\exp(C_hn_e)\). It is independent of \(\eta\) and \(M\).

At \(\eta=0\), part 1 of
Proposition~\ref{prop:04-integer-bounds} gives
\begin{equation}
 |F_{m'(M)}(M)|
 \le\exp\!\left(C_hL+C_hn_e
       -c_h\left(\frac{M^2+n_e^2}{L}+\ell_e\right)\right).
 \label{eq:06-real-sector}
\end{equation}
Its dimension range \(0\le m'_\sigma\le C e^{\gamma L}\) includes
every valid sector because \(e^{\gamma L}=N\). In particular,
\eqref{eq:06-real-sector} requires no polynomial restriction on
\(M\), \(n_e\), or \(\ell_e\). That assertion uses the coarse real
reference-source inequality and compensated quadratic integration;
the later sharp source estimate and interpolation of dimensions are
not needed here.

The phase of the exact sector identity is
\(\exp(i(\eta+\pi f)M)\), where the far-particle integer \(f\) is
fixed by \(\zeta\). Replacing \(\eta=0\) by \(iv\) therefore
multiplies its modulus by \(e^{-vM}\le e^{v|M|}\). Every other factor
in \eqref{eq:06-real-sector} is unchanged. Combining the prefactor
and this estimate bounds the divided numerator by
\[
 e^{C_hL}
 \sum_{M\in\mathbb Z}e^{-c_hM^2/L+v|M|}
 \sum_{\zeta}\!\int
       e^{C'_hn_e-c_h\ell_e}.
\]
We have discarded the nonpositive \(n_e^2\) term and relaxed the
valid-count constraints, which only enlarges the positive majorant.
All integrals in the last sum retain their species factorials.

Here is the explicit-configuration estimate. For a far species its
real parameter ranges over \(\mathbb R\), and
\[
 \int_{\mathbb R}e^{-c_h(|u|-L)_+}\,du=2L+2/c_h.
\]
For an exterior main species it ranges over \(|u|>L\), giving
\(2/c_h\). Fixed contour-measure and activity constants can be
absorbed into \(e^{C'_h}\). For each of the finitely many species,
with parameter domain \(D_I\), the factorial sum is at most
\[
 \sum_{n=0}^{\infty}\frac{e^{C'_hn}}{n!}
 \left(\int_{D_I}e^{-c_h(|u|-L)_+}\,du\right)^n
 \le e^{C''_h(L+1)}.
\]
Taking the product gives \(e^{C_hL}\), for \(L\ge1\).
The elementary multiplicity two of a fused double affects the
admissibility constraint already dropped; it does not remove its
factorial divisor or introduce a second integration variable.

Only the imbalance sum remains. Young's inequality gives
\[
 -c_hM^2/L+v|M|
 \le-\frac{c_hM^2}{2L}+\frac{Lv^2}{2c_h},
\]
so
\[
 \sum_{M\in\mathbb Z}e^{-c_hM^2/L+v|M|}
 \le C_h\sqrt L\,e^{C_hLv^2}.
\]
Absorbing \(\sqrt L\) into \(e^{C_hL}\) proves
\( |G_N(iv)|\le e^{C_hL(1+v^2)}\) for the divided numerator
\(G_N\) in \eqref{eq:G-22}, uniformly for all \(v\ge0\).

Finally divide by the \(\eta\)-independent denominator \(\Delta_N\).
The balanced denominator theorem, Theorem~\ref{thm:02-denominator},
gives \(\log|\Delta_N|=(5/24)\log N+o(\log N)\). In particular
it is nonzero and \(|\Delta_N|^{-1}\le e^{C\log N}\) for all
sufficiently large even \(N\). This bound is independent of \(v\).
The physical value \(Z_N(2\cosh v)\) is nonnegative by
\eqref{eq:06-positive-coefficients}. Dividing the numerator estimate
and using \(L=(\log N)/\gamma\) proves
\eqref{eq:06-all-fugacity}.
\end{proof}

\begin{corollary}[Tail in the polygon count]\label{E:all-count-tail}
\label{lem:11-polygon-tail}
There are constants \(c,C>0\) such that for every sufficiently large
even \(N\) and every integer \(a\ge0\),
\begin{equation}
 w_N(a)\le
 \exp\!\left(C\log N-\frac{ca^2}{\log N}\right).
 \label{eq:L-40}
\end{equation}
\end{corollary}
\begin{proof}
For every \(v\ge0\), positivity in
\eqref{eq:06-positive-coefficients} gives
\[
 w_N(a)(2\cosh v)^a\le Z_N(2\cosh v).
\]
Since \(\log(2\cosh v)\ge v\),
Theorem~\ref{E:all-fugacity-bound} yields
\[
 w_N(a)\le\exp\!\left(C\log N(1+v^2)-av\right).
\]
Take \(v=a/(2C\log N)\). The exponent becomes
\(C\log N-a^2/(4C\log N)\), as required. This includes \(a=0\);
indices beyond the polynomial degree have zero mass.
\end{proof}
